\documentclass[aspectratio=169,10pt]{beamer}

% ============================================================
%  PAQUETES
% ============================================================
\usepackage[T1]{fontenc}
\usepackage[utf8]{inputenc}
\usepackage{amsmath,amssymb,bm}
\usepackage{tikz}
\usetikzlibrary{arrows.meta,positioning,calc,backgrounds,shapes.geometric}
\usepackage{tcolorbox}
\tcbuselibrary{skins,breakable}
\usepackage{booktabs}
\usepackage{colortbl}
\usepackage{listings}
\usepackage{xcolor}

% ============================================================
%  PALETA
% ============================================================
\definecolor{darkbg}{HTML}{1A1A2E}
\definecolor{accent}{HTML}{E94560}
\definecolor{accentteal}{HTML}{0F9B8E}
\definecolor{accentamber}{HTML}{F5A623}
\definecolor{accentviolet}{HTML}{7B68EE}
\definecolor{textlight}{HTML}{EAEAEA}
\definecolor{textgray}{HTML}{9E9E9E}
\definecolor{lightbg}{HTML}{F7F7F7}
\definecolor{coralbox}{HTML}{FF6B6B}
\definecolor{codebg}{HTML}{272740}
\definecolor{codefg}{HTML}{E8E8F0}

% ============================================================
%  TEMA BEAMER
% ============================================================
\setbeamercolor{background canvas}{bg=lightbg}
\setbeamercolor{normal text}{fg=darkbg}
\setbeamertemplate{navigation symbols}{}
\setbeamertemplate{footline}{%
	\hfill\textcolor{textgray}{\tiny\texttt{mbastidaso@unal.edu.co}%
		\quad\insertframenumber/\inserttotalframenumber\quad}%
	\vspace{4pt}%
}
\setbeamertemplate{itemize item}{\textcolor{accent}{\textbullet}}
\setbeamertemplate{itemize subitem}{\textcolor{accentteal}{--}}

% ============================================================
%  ENTORNO DARK SLIDE
% ============================================================
\newenvironment{darkslide}{%
	\setbeamercolor{background canvas}{bg=darkbg}%
	\setbeamercolor{normal text}{fg=textlight}%
}{%
	\setbeamercolor{background canvas}{bg=lightbg}%
	\setbeamercolor{normal text}{fg=darkbg}%
}

% ============================================================
%  CAJAS DE CONTENIDO
% ============================================================
\tcbset{mybase/.style={
		enhanced, breakable, arc=3pt,
		boxrule=0.6pt, left=6pt, right=6pt, top=5pt, bottom=5pt,
		fonttitle=\bfseries, fontupper=\normalsize
}}
\newtcolorbox{defbox}[1]{mybase,
	colback=accentviolet!8,colframe=accentviolet!60,title={#1}}
\newtcolorbox{resultbox}[1]{mybase,
	colback=accentteal!8,colframe=accentteal!60,title={#1}}
\newtcolorbox{obsbox}[1]{mybase,
	colback=accentamber!10,colframe=accentamber!70,title={#1}}
\newtcolorbox{taskbox}[1]{mybase,
	colback=coralbox!8,colframe=coralbox!60,title={#1}}

% ============================================================
%  CÓDIGO
% ============================================================
\lstset{
	language=Python,
	basicstyle=\ttfamily\scriptsize\color{codefg},
	backgroundcolor=\color{codebg},
	keywordstyle=\color{accentviolet!90!white}\bfseries,
	stringstyle=\color{accentteal},
	commentstyle=\color{textgray}\itshape,
	frame=none, numbers=none,
	breaklines=true, showstringspaces=false,
	xleftmargin=4pt, xrightmargin=4pt,
	aboveskip=3pt, belowskip=3pt,
	morekeywords={np,plt,self}
}

\newcommand{\R}{\mathbb{R}}

% ============================================================
\begin{document}
	% ============================================================
	
	% ── PORTADA ──────────────────────────────────────────────────
	\begin{darkslide}
		\begin{frame}[plain]
			\begin{tikzpicture}[remember picture, overlay]
				\fill[accent,opacity=0.15]      (current page.north east) ++(-1.5,-1.5) circle (4cm);
				\fill[accentteal,opacity=0.10]  (current page.south west) ++(2,2)       circle (3cm);
				\fill[accentviolet,opacity=0.08](current page.north west) ++(1,-2)       circle (2.5cm);
			\end{tikzpicture}
			\vspace{1.5cm}
			\begin{center}
				{\small\textcolor{textgray}{Programaci\'on Cient\'ifica $\cdot$ 2026-1}}\\[16pt]
				{\LARGE\textcolor{textlight}{\textbf{Repaso: EDOs, el chisme y POO}}}\\[6pt]
				{\large\textcolor{accent}{M\'odulo 4 -- Sesi\'on integradora}}\\[14pt]
				{\normalsize\textcolor{textgray}{%
						De la ecuaci\'on al c\'odigo, y del c\'odigo a la clase}}\\[20pt]
				{\normalsize\textcolor{textgray}{Universidad Nacional de Colombia}}\\[4pt]
				{\small\textcolor{textgray}{\texttt{mbastidaso@unal.edu.co}}}
			\end{center}
		\end{frame}
	\end{darkslide}
	
	% ── PARTE 1: REPASO CONCEPTUAL ───────────────────────────
	\begin{darkslide}
		\begin{frame}[plain]
			\begin{tikzpicture}[remember picture, overlay]
				\fill[accentteal,opacity=0.12](current page.south east)++(-2,2)  circle(3.5cm);
				\fill[accent,opacity=0.08]    (current page.north west)++(2,-1)  circle(2.5cm);
			\end{tikzpicture}
			\vspace{2.2cm}
			\begin{center}
				{\Large\textcolor{textlight}{\textbf{Parte 1}}}\\[8pt]
				{\large\textcolor{accentteal}{Repaso: EDOs y sistemas din\'amicos}}
			\end{center}
		\end{frame}
	\end{darkslide}
	
	% ─────────────────────────────────────────────────────────
	\begin{frame}{Recordatorio r\'apido}
		\vspace{4pt}
		\begin{defbox}{Lo que ya vimos}
			Un sistema din\'amico describe c\'omo cambia un \textbf{estado}
			$x(t)$ en el tiempo mediante una EDO:
			\[
			\dot{x}(t) = f\!\bigl(x(t),t\bigr), \qquad x(t_0)=x_0.
			\]
			Si $f$ es Lipschitz, Picard-Lindel\"of garantiza soluci\'on \'unica.
		\end{defbox}
		
		\vspace{8pt}
		\begin{taskbox}{Hoy el estado es distinto}
			En vez de una opini\'on continua (H-K), cada persona est\'a en uno de
			\textbf{tres estados discretos} frente a un chisme: no lo sabe, lo
			cuenta, o ya se cans\'o de contarlo.
		\end{taskbox}
	\end{frame}
	
	% ─────────────────────────────────────────────────────────
	\begin{frame}{El chisme como sistema din\'amico}
		\vspace{4pt}
		\textbf{Tres estados (fracciones de la poblaci\'on):}
		\begin{itemize}\small
			\item \textbf{Ignorante} ($i$): no conoce el chisme.
			\item \textbf{Difusor} ($s$): lo conoce y lo cuenta.
			\item \textbf{Cansado} ($r$): lo conoce pero ya no lo cuenta.
		\end{itemize}
		
		\vspace{6pt}
		\textbf{Reglas (Maki-Thompson):}
		\begin{itemize}\small
			\item Ignorante $+$ difusor $\to$ dos difusores (tasa $\beta$).
			\item Difusor $+$ alguien informado $\to$ un cansado m\'as
			(tasa $\alpha$): se cansa al notar que ya lo sab\'ian.
		\end{itemize}
		
		\vspace{6pt}
		\begin{defbox}{Modelo}
			\[
			\dot{i} = -\beta i s, \qquad
			\dot{s} = \beta i s - \alpha s(s+r), \qquad
			\dot{r} = \alpha s(s+r), \qquad i+s+r=1.
			\]
		\end{defbox}
	\end{frame}
	
	% ─────────────────────────────────────────────────────────
	\begin{frame}{Repaso: Euler y RK4}
		\vspace{4pt}
		\begin{columns}[T]
			\column{0.48\textwidth}
			\begin{resultbox}{Euler (orden 1)}
				\begin{align*}
					w_0   &= x_0, \\
					w_{i+1} &= w_i + h\,f(t_i,w_i)
				\end{align*}
				Error global $O(h)$.
			\end{resultbox}
			
			\column{0.48\textwidth}
			\begin{resultbox}{RK4 (orden 4)}
				\begin{align*}
					w_0   &= x_0, \\
					k_1   &= h\,f(t_i,\;w_i), \\
					k_2   &= h\,f\!\left(t_i + \tfrac{h}{2},\; w_i + \tfrac{1}{2}k_1\right), \\
					k_3   &= h\,f\!\left(t_i + \tfrac{h}{2},\; w_i + \tfrac{1}{2}k_2\right), \\
					k_4   &= h\,f\!\left(t_i + h,\;             w_i + k_3\right), \\
					w_{i+1} &= w_i + \tfrac{1}{6}\!\left(k_1 + 2k_2 + 2k_3 + k_4\right),
				\end{align*}
				Error global $O(h^4)$.
			\end{resultbox}
		\end{columns}
		
	\end{frame}
	
	% ── PARTE 3: DE FUNCIONES A CLASES ───────────────────────
	\begin{darkslide}
		\begin{frame}[plain]
			\begin{tikzpicture}[remember picture, overlay]
				\fill[coralbox,opacity=0.12]    (current page.north east)++(-2,-2) circle(3.5cm);
				\fill[accentviolet,opacity=0.10](current page.south west)++(1,1)   circle(3cm);
			\end{tikzpicture}
			\vspace{2cm}
			\begin{center}
				{\Large\textcolor{textlight}{\textbf{Parte 2}}}\\[8pt]
				{\large\textcolor{coralbox}{De funciones a clases}}
			\end{center}
		\end{frame}
	\end{darkslide}
	
	% ─────────────────────────────────────────────────────────
	\begin{frame}[fragile]{El patr\'on que se repite}
		
		\vspace{2pt}
		En el notebook, \texttt{euler} y \texttt{rk4} son funciones sueltas
		que reciben siempre lo mismo:
		
		\begin{lstlisting}
			def euler(f, x0, t, **kw): ...
			def rk4(f, x0, t, **kw): ...
		\end{lstlisting}
		
		\vspace{6pt}
		\begin{obsbox}{Idea}
			Si el campo \texttt{f}, el estado y el paso \texttt{dt} siempre viajan
			juntos, agrup\'emoslos en un objeto que \textbf{recuerda su propio
				estado} entre llamadas.
		\end{obsbox}
	\end{frame}
	
	% ─────────────────────────────────────────────────────────
	\begin{frame}[fragile]{Ingredientes m\'inimos de una clase}
		\vspace{2pt}
		\begin{columns}[T]
			\column{0.46\textwidth}
			\begin{defbox}{Vocabulario}
				\begin{itemize}\small
					\item \texttt{class}: define un nuevo tipo.
					\item \texttt{\_\_init\_\_}: se ejecuta al crear el objeto.
					\item \texttt{self}: el objeto mismo.
					\item \textbf{atributo}: dato que vive en el objeto.
					\item \textbf{m\'etodo}: funci\'on que vive en el objeto.
				\end{itemize}
			\end{defbox}
			
			\column{0.50\textwidth}
\begin{lstlisting}
class Particula:
def __init__(self, x, v):
self.x = x
self.v = v

def mover(self, dt):
self.x += self.v * dt

p = Particula(x=0.0, v=2.0)
p.mover(dt=0.5)
print(p.x)   # 1.0
\end{lstlisting}
		\end{columns}
	\end{frame}
	
	% ─────────────────────────────────────────────────────────
	\begin{frame}[fragile]{La clase \texttt{Solver} y su subclase \texttt{Euler}}
		\vspace{2pt}
\begin{lstlisting}
class Solver:
def __init__(self, f, x0, dt, **kwargs):
self.f, self.dt, self.kwargs = f, dt, kwargs
self.x = np.array(x0, dtype=float)
self.t = 0.0

def step(self):
raise NotImplementedError(
"Las subclases deben implementar step()")


class Euler(Solver):
def step(self):
self.x = self.x + self.dt * self.f(
self.x, self.t, **self.kwargs)
self.t += self.dt
\end{lstlisting}
		
		\vspace{4pt}
		\begin{resultbox}{Herencia}
			\texttt{Euler} hereda \texttt{\_\_init\_\_} sin reescribirlo y solo
			sobrescribe \texttt{step()}.
		\end{resultbox}
	\end{frame}
	
	% ─────────────────────────────────────────────────────────
	\begin{frame}{Ejercicio: complete \texttt{RK4(Solver)}}
		\vspace{4pt}
		\begin{taskbox}{En el notebook}
			Complete \texttt{step()} de \texttt{RK4(Solver)} con la f\'ormula
			de RK4, reutilizando \texttt{self.f}, \texttt{self.x}, \texttt{self.t},
			\texttt{self.dt}, \texttt{self.kwargs}. Verifique que coincide con
			\texttt{rk4()} aplicada al mismo \texttt{dk\_rhs} que ya escribieron.
		\end{taskbox}
		
		\vspace{10pt}
		\begin{obsbox}{El punto de toda la sesi\'on}
			El mismo \texttt{Solver} que integr\'o el chisme sirve para
			\textbf{cualquier} modelo futuro solo cambiando \texttt{f}.
		\end{obsbox}
	\end{frame}
	
	% ── SLIDE FINAL ───────────────────────────────────────────
	\begin{darkslide}
		\begin{frame}[plain]
			\begin{tikzpicture}[remember picture, overlay]
				\fill[accent,opacity=0.15]      (current page.north east) ++(-1,-1) circle (3.5cm);
				\fill[accentteal,opacity=0.10]  (current page.south west) ++(2,2)   circle (2.8cm);
				\fill[accentviolet,opacity=0.08](current page.north west) ++(1,-2)  circle (2cm);
			\end{tikzpicture}
			\vspace{1.8cm}
			\begin{center}
				{\large\textcolor{textgray}{Recapitulando}}\\[14pt]
				{\large\textcolor{textlight}{\textbf{Modelo del chisme}}}\\
				{\small\textcolor{textgray}{tres estados discretos, campo suave, un rumor nunca llega al 100\%}}\\[10pt]
				{\large\textcolor{textlight}{\textbf{Clases y herencia}}}\\
				{\small\textcolor{textgray}{\texttt{Solver} agrupa estado y comportamiento; las subclases sobrescriben \texttt{step()}}}\\[10pt]
				{\large\textcolor{textlight}{\textbf{Un patr\'on, muchos modelos}}}\\
				{\small\textcolor{textgray}{el mismo c\'odigo servir\'a para todo lo que viene en el curso}}\\[20pt]
				\rule{4cm}{0.4pt}\\[6pt]
				{\small\textcolor{textgray}{\texttt{mbastidaso@unal.edu.co}}}
			\end{center}
		\end{frame}
	\end{darkslide}
	
\end{document}